\clearpage
\chapter*{Errata da exposição n\textsuperscript{o} 190}
\addcontentsline{toc}{chapter}{Errata da exposição n\textsuperscript{o} 190}
\setcounter{footnote}{0}
\markboth{Errata da exposição 190}{Errata da exposição 190}
\sid{e190}
\src{300}

\noindent\textbf{3. Técnica de descida e teoremas de existência em geometria
algébrica [TDTE I]. I: Generalidades. Descida por morfismos fielmente planos.}

\begin{center}
[Seminário Bourbaki, t. 12, 1959/60, n\textsuperscript{o} 190, 29 p.]
\end{center}

\noindent\textbf{Página 1, linha 10.} --- Agora parece excessivo dizer que a
técnica de descida está «na base da maioria dos teoremas de
existência em geometria algébrica». Isto é verdadeiro em grande medida para as
técnicas não projetivas tratadas nas duas primeiras exposições
da série TDTE I a VI, mas não para as técnicas projetivas (TDTE IV, V,
VI).

\noindent\textbf{Página 5, linha 16.} --- Não é necessário supor que $\alpha$
seja um morfismo de $\mathcal F$-descida.

\noindent\textbf{Página 20, observação.} --- Um morfismo
$S'\longrightarrow S$ quase-finito, étale e sobrejetivo, ou um morfismo
$\operatorname{Spec}(\widehat A)\longrightarrow\operatorname{Spec}(A)$, não é
sempre um morfismo de descida estrito, mesmo que $A$ seja um anel local de uma
curva algébrica sobre um corpo algebraicamente fechado $k$ e
$S=\operatorname{Spec}(A)$. Assim, pode-se encontrar um morfismo próprio e liso
$f:X\longrightarrow S$, que faz de $X$ um fibrado principal de base $S$ sob
uma curva elíptica $E$, tal que $f':X'\longrightarrow S'$ seja projetivo, mas
$f$ não o seja. É, portanto, também um exemplo de fibrado
principal homogêneo não isotrivial sob um esquema abeliano.

\noindent\textbf{Página 26, linha 9.} --- Leia-se «Chow--Lang» em vez de
«Igusa--Lang».

Para diversos detalhes acerca da teoria da descida, vejam-se SGA VI, VII e
VIII.

\par
